\section{Mechanism 1 in the cube worlds: the state is kept only while it is predictable}
\label{sec:representation}

\paragraph{Design.}
Four gray-noise conditions cross background persistence (per-frame versus persistent) with
dynamics (deterministic versus hidden wind). Each condition receives three independently
initialized grid JEPAs with identical architecture, 512-dimensional latent, 500,000 transitions
and 8,500 updates. After training, a nonlinear position read-out is fitted on frozen encodings of
training frames and scored on held-out frames as L1 error in pixels; the same read-out is then
applied to eight-step predicted latents (Appendix~\ref{app:cubes}). Nothing about the read-out
depends on the condition, so a difference between cells is a difference in what the encoder
exposes.

\begin{table}[t]
\centering
\caption{\textbf{The effect of nuisance persistence reverses under stochastic dynamics.} Cube
position L1 error in pixels, mean $\pm$ sample standard deviation over three model fits. ``Real''
reads encodings of held-out observations; the right column applies the same read-out eight
predicted steps ahead. Architecture, latent size, data and updates are fixed across cells.}
\label{tab:cube-main}
\inputtable{cube_main.tex}
\end{table}

\paragraph{The reversal.}
With deterministic motion the objective keeps the cube in every condition, and per-frame noise is
the \emph{easier} background: $0.115\pm0.012$ pixels against $0.764\pm0.068$ for persistent noise
(Table~\ref{tab:cube-main}, Figure~\ref{fig:cube-readout}). Adding wind flips the ordering and
breaks one cell outright. Under persistent noise the read-out still localizes the cube to
$0.668\pm0.049$ pixels; under per-frame noise it fails at $21.33\pm0.06$ pixels, an error larger
than the cube is wide in every seed, with almost no seed-to-seed spread. The failure is already
present in the encodings of real frames, before any prediction is made: decoded real encodings from
a failing model show a gray field with at most a faint trace of the cube, and its predicted future
is not a scene at all, while the flow model trained on the same data keeps the cube in both
(Figure~\ref{fig:wind-filmstrip}). The eight-step prediction column tells the same story one step
later: where the state was kept, prediction error grows with the horizon as the wind's uncertainty
accumulates; where it was dropped, prediction has nothing left to lose.

\begin{figure}[t]
\centering
\paperfig{gray_frame_wind_filmstrip.png}{1.0}
\caption{\textbf{Per-frame noise plus wind: the JEPA no longer encodes the cube.} Recorded future
under the recorded actions; the JEPA's decoded encodings of those real frames; the JEPA's decoded
rollout; the flow model's decoded encodings and one sampled future. First eligible episode of the
held-out set, chosen without looking at any model output; JEPA model seed 0, flow seed 0. The
decoders are supervised reconstruction probes fitted after training, so the middle rows show what
is recoverable from the frozen latent, not what the model was trained to produce.}
\label{fig:wind-filmstrip}
\end{figure}

\begin{figure}[t]
\centering
\paperfig{imagined_termination.pdf}{1.0}
\caption{\textbf{What each model's imagination tells a policy about dying.} Each frozen model is
rolled forward from 256 real contexts under the recorded (uniformly random) actions and its own
termination head is read on its own predicted latents; the curve is the implied probability that the
imagined episode is still alive after $k$ steps, against the recorded survival of the same episodes
(dashed). In the per-frame-noise wind world the JEPA head reads a $0.91$ termination probability on
every predicted latent from the first imagined step, while on encoder latents of the real frames it
tracks the recorded rate ($0.03$--$0.06$); the model has no state left for the head to read, so
imagination ends in two steps whatever the action. Under persistent noise the JEPA keeps the cube
and its deterministic rollout is optimistic instead: the hidden wind's excursions are averaged away
(Appendix~\ref{app:shared-cnn}).}
\label{fig:imagined-termination}
\end{figure}

\paragraph{Why this is mechanism 1.}
Under deterministic motion the next position is a function of the history and the action, so
keeping it makes the target predictable and lowers the loss. Wind removes that incentive: the
position's innovations cost prediction loss no matter how the encoder represents them
(Corollary~\ref{cor:strict}). The failing cell is not the zero-loss nuisance optimum of
Proposition~\ref{prop:representation}: its prediction loss stays high, about $0.733$ per
coordinate, SIGReg is satisfied and the latent is far from degenerate. The objective did not find a
cheap substitute for the state; it simply stopped paying for the state once paying no longer bought
predictability. Why persistent noise protects the cube under the same wind is not something the
theory decides. The matched design establishes only that it does, in all three seeds, and that
the protection is an interaction between nuisance and dynamics rather than a property of either.

\paragraph{The state is encoded through the target.}
A second intervention isolates where the incentive acts. In the per-frame conditions we refit
three seeds per world with the prediction target detached from the encoder's gradient, changing
nothing else (the predictor-input and SIGReg gradients remain). With deterministic motion the
detached models lose most of the position information that the joint models keep; with wind, both
variants fail identically (Table~\ref{tab:target-gradient}). The encoder keeps the cube only when
the joint gradient can make the cube's future predictable, and the wind failure is not a side
effect of that gradient path. Detaching the target, the natural first repair, does nothing here.

\begin{table}[t]
\centering
\caption{\textbf{Detaching the target does not rescue the wind cell and costs the deterministic
one.} Position read-out error (mean absolute pixel error per coordinate, read-out fitted on
50,000 frames; a different read-out protocol from Table~\ref{tab:cube-main}), mean $\pm$ sample
standard deviation over three model seeds; each seed pair shares initialization and data order.}
\label{tab:target-gradient}
\inputtable{target_gradient.tex}
\end{table}

\paragraph{Behavioral consequence.}
A policy that cannot see the cube cannot steer it. \pendingresult{Insert the shared-CNN outcomes
for the three cube worlds (per-frame still, per-frame wind, persistent wind) from the running
campaign: survival and control length for JEPA and flow imagination, three policy seeds each,
paired keys. State the interface qualification. Appendix~\ref{app:cubes} keeps the single-run
survival audit across all six cube worlds.}
